% ============================================================
% A26 — Scheduler dan Queue
% ============================================================
\flowmeta{A26}{Scheduler dan Queue}{Admin/Ops}

\begin{flowsection}{Tujuan}
Menjalankan tugas otomatis terjadwal dan antrean (queue) agar proses latar
berjalan tanpa intervensi manual.
\end{flowsection}

\begin{flowsection}{Prasyarat}
Cron server aktif memanggil \texttt{schedule:run} dan worker queue berjalan.
\end{flowsection}

\begin{flowsection}{Tugas Terjadwal}
\begin{itemize}
  \item \texttt{orders:cancel-expired} --- membatalkan pesanan kedaluwarsa.
  \item \texttt{loyalty:expire-points} --- memproses poin kedaluwarsa.
  \item \texttt{shipments:sync-tracking} --- menyinkronkan status pelacakan.
  \item \texttt{expeditions:sync} --- menyinkronkan data ekspedisi/kurir.
  \item \texttt{queue:work} --- memproses antrean (mis. pengiriman email).
  \item \texttt{model:prune} --- membersihkan data sementara.
\end{itemize}
\end{flowsection}

\shot{gambar/admin/a26-1.png}{Diagram penjadwalan tugas otomatis (scheduler) dan antrean (queue)}{fig:a26-1}

\begin{flowsection}{Hasil yang Diharapkan}
Tugas berjalan terjadwal; antrean email/webhook terproses tanpa memblokir permintaan pengguna.
\end{flowsection}

\begin{flowsection}{Penanganan Error dan Kasus Khusus}
\begin{itemize}
  \item Cron tidak berjalan $\rightarrow$ tugas tidak dieksekusi (cek crontab server).
  \item Queue menumpuk $\rightarrow$ pastikan worker aktif; gunakan \texttt{withoutOverlapping}.
\end{itemize}
\end{flowsection}

\begin{flowsection}{Kaitan Modul}
\texttt{backend/routes/console.php}, \texttt{backend/app/Console/Commands/*},
cron \texttt{schedule:run}.
\end{flowsection}

\docver{v1.0}{Sep 2026}
